% ═══════════════════════════════════════════════════════════════════════
% IntLLM Paper — MLSys 2027 / arXiv submission
%
% Per FJQ_PHASE_E_PATH_A_PAPER_OUTLINE.md v1.0 §3:
%   10-page MLSys + 11-page arXiv submission
%   Five contributions C1-C5
%   Submission target: 2026-05-25
%
% Status (2026-04-27 — Path A Week 1, sub-step 2):
%   §3 architecture     — DRAFTED THIS COMMIT
%   §4 methodology      — STUB (next sub-step)
%   §1 §2 §5-§10        — STUB (Week 1-3 calendar per outline §7)
%
% This file is the manuscript's LaTeX source. Tables are auto-generated
% from JSON artifacts via `paper/intllm/tables/*.md` + verify_intllm_tables.py;
% each numeric literal in this file MUST correspond to a registered Claim
% per the §6.9 R7 audit gate.
% ═══════════════════════════════════════════════════════════════════════

\documentclass[10pt,twocolumn]{article}

% ── Packages ──
\usepackage[utf8]{inputenc}
\usepackage[T1]{fontenc}
\usepackage{amsmath,amssymb,amsthm}
\usepackage{booktabs}
\usepackage{hyperref}
\usepackage{graphicx}
\usepackage{natbib}
\usepackage[margin=1in]{geometry}
\usepackage{enumitem}
\usepackage{listings}
\usepackage{makecell}
\usepackage{multirow}
\usepackage{balance}

% ── Code listings (consistent with v3.1 paper) ──
\lstset{basicstyle=\ttfamily\footnotesize,breaklines=true,frame=single}

% ── Theorem environments (mostly unused; kept for parity with v3.1) ──
\newtheorem{theorem}{Theorem}
\newtheorem{lemma}[theorem]{Lemma}
\newtheorem{definition}[theorem]{Definition}

% ── Title ──
\title{IntLLM: Compiler-Verified 1.58-Bit Ternary Language Models for\\
       Operating-System Kernel Inference}
\author{Muhamad Fajar Putranto\\TaxPrime / PrimeCore.id, Jakarta, Indonesia}
\date{}

\begin{document}
\maketitle

% ═══════════════════════════════════════════════════════════════════════
% Abstract (STUB — drafted Week 3 per Path A outline §7)
% ═══════════════════════════════════════════════════════════════════════

\begin{abstract}
\textit{[Abstract pending — drafted Week 3 of Path A calendar after the
results sections (§5-§7) land. Will summarize the five contributions:
(C1) Phase D scaling chain at training-time-quant validates
Chinchilla-style behavior at 22M-74M parameter range; (C2) bilingual
ID+EN extension via Phase E1 corpus v1.0 (25.67~B tokens at 60:40);
(C3) in-kernel deployment via FajarOS Nova v3.9.0 IntLLM kernel-path;
(C4) honest negative results on outlier-handling features
(\texttt{balanced\_calib} and Hadamard rotation) at this regime, with
3-cause attribution; (C5) compiler-verified safety via
\texttt{@kernel}/\texttt{@device}/\texttt{@noinline} annotations.]}
\end{abstract}

% ═══════════════════════════════════════════════════════════════════════
\section{Introduction}
\label{sec:intro}

\textit{[STUB — drafted Week 3 of Path A calendar. Motivates the
kernel-resident LLM use case + explains why existing solutions
(BitNet~2B4T, MMfreeLM at $\geq$370M, eBPF-AI userspace offload) do
not meet the $\leq$16~MB binary + zero-heap +
deterministic-execution constraints simultaneously. Sets up C1-C5
contribution preview.]}

% ═══════════════════════════════════════════════════════════════════════
\section{Background}
\label{sec:background}

\textit{[STUB — drafted Week 3 of Path A calendar. Covers BitNet
\citep{wang2023bitnet}, BitNet~b1.58~2B4T \citep{ma2024era},
MatMul-Free LLM \citep{zhu2024matmulfree}, and the
kernel-context constraints (no heap allocation during inference,
$\leq$16~MB total binary footprint, deterministic
execution, Fajar Lang \texttt{@kernel} annotation system).]}

% ═══════════════════════════════════════════════════════════════════════
% §3 IntLLM Architecture — DRAFTED THIS COMMIT (Path A Week 1 sub-step 2)
% ═══════════════════════════════════════════════════════════════════════
\section{IntLLM Architecture}
\label{sec:arch}

IntLLM is a decoder-only language model family that extends the
MatMul-Free LLM architecture \citep{zhu2024matmulfree} with full
integer-native operation, kernel-context-friendly memory layout, and a
calibrated ternary quantization pipeline that runs end-to-end inside
the FajarOS Nova \texttt{@kernel} address space.

\subsection{Model class hierarchy}
\label{sec:arch-class}

The training-time class hierarchy is
\texttt{HGRNBitForCausalLM}~$\to$~\texttt{HGRNBitModel}~$\to$~\texttt{HGRNBitBlock},
inherited verbatim from the upstream MatMul-Free LLM repository (frozen
pin in \texttt{python/phase\_d/\_upstream/mmfreelm/}). IntLLM extends
the upstream surface in three places without modifying upstream
sources: (i)~a custom training driver
\texttt{intllm.train.train\_loop} that adds calibrated gates per scale,
checkpoint atomicity, and a step-idle watchdog (per the cross-repo
production-safety rule of Section~\ref{sec:method-trackb}); (ii)~a
quantization helper module \texttt{intllm.quant} adding RMSNorm
variants, sigmoid/SiLU lookup tables, and the per-channel adaptive
quantizer (Section~\ref{sec:arch-bitlinear}); and (iii)~an evaluation
harness \texttt{intllm.eval} aligning measurements with
\texttt{lm-evaluation-harness~v0.4.11} for baseline parity per the
research-integrity rules of Section~\ref{sec:method-integrity}.

\subsection{BitLinear forward path}
\label{sec:arch-bitlinear}

The fundamental quantized operator is
\texttt{BitLinear}~\citep{wang2023bitnet,ma2024era}, a drop-in
replacement for \texttt{nn.Linear} that quantizes both weights and
activations during training. The forward path is

\begin{equation}
\label{eq:bitlinear}
y = \text{F.linear}\bigl(\widehat{x}_{\text{quant}},\;\widehat{W}_{\text{quant}}\bigr),
\end{equation}

where the activation and weight quantizers, with straight-through
estimator (STE) gradient passthrough, are

\begin{align}
x_{\text{norm}} &= \text{RMSNorm}(x), \\
\widehat{x}_{\text{quant}} &= x_{\text{norm}} + \text{StopGrad}\bigl(q_a(x_{\text{norm}}) - x_{\text{norm}}\bigr), \\
\widehat{W}_{\text{quant}} &= W + \text{StopGrad}\bigl(q_w(W) - W\bigr).
\end{align}

The activation quantizer $q_a$ is per-token absmax 8-bit:
\begin{equation}
\label{eq:actquant}
q_a(x)_{t,c} = \frac{\bigl\lfloor x_{t,c}\cdot s_t \bigr\rceil_{[-128,127]}}{s_t},
\quad
s_t = \frac{127}{\max_{c'} |x_{t,c'}| + \epsilon},
\end{equation}
where $\lfloor\cdot\rceil_{[-128,127]}$ denotes round-then-clamp.
The weight quantizer $q_w$ is per-tensor absmean ternary
$\{-1, 0, +1\}$ per b1.58:
\begin{equation}
\label{eq:wquant}
q_w(W)_{i,j} = \frac{\bigl\lfloor W_{i,j}\cdot s_w \bigr\rceil_{[-1,1]}}{s_w},
\quad
s_w = \frac{1}{\overline{|W|} + \epsilon},
\end{equation}
giving an effective $\log_2 3 \approx 1.58$ bits per weight after entropy
coding. We use $\epsilon = 10^{-5}$ throughout, matching the upstream
implementation.

\subsection{Per-block layout}
\label{sec:arch-block}

Each \texttt{HGRNBitBlock} contains exactly six BitLinear instances:
four in the attention sub-layer (\texttt{i\_proj}, \texttt{f\_proj},
\texttt{g\_proj}, \texttt{o\_proj}) implementing the gated linear
recurrence of MMfreeLM~\citep{zhu2024matmulfree}, and two in the
feed-forward sub-layer (\texttt{mlp.gate\_proj},
\texttt{mlp.down\_proj}). A top-level \texttt{lm\_head} BitLinear
projects to the vocabulary, giving a per-model BitLinear count of
$6L + 1$, where $L$ is the number of layers. For our Mini configuration
($L = 6$), this is $37$ BitLinear sites; this count is recoverable as a
runtime invariant via \texttt{intllm.qat.is\_bitlinear}-based traversal
and is reported by every Phase~D ablation harness.

The structural distinction between \texttt{i\_proj}/\texttt{f\_proj}/\texttt{g\_proj}
and \texttt{o\_proj} is load-bearing for our negative-result analysis
in Section~\ref{sec:negatives}: the first three implement gated linear
recurrence projections (structurally distinct from the QKV projections
of softmax attention) while \texttt{o\_proj} is the post-recurrence
output projection. Outlier-handling rotations from the QuaRot family
\citep{ashkboos2024quarot,liu2024spinquant} are canonically applied
to QKV inputs in transformers; we test the closest analogue (\texttt{o\_proj}
input rotation) at \mbox{Mini scale} and find it insufficient
(Section~\ref{sec:negatives}).

\subsection{Tokenizer and vocabulary}
\label{sec:arch-tokenizer}

We use the Mistral~v3 byte-pair encoding tokenizer with 32{,}768 token
vocabulary, chosen because (i)~it has wide multilingual coverage at a
size that fits the kernel-binary budget, (ii)~it is the same tokenizer
used by \citet{zhu2024matmulfree} so cross-paper benchmarks are
directly comparable, and (iii)~its training data spans Latin-script
languages including Indonesian, supporting the bilingual extension of
Section~\ref{sec:bilingual} without retraining the tokenizer.

\subsection{Configuration matrix}
\label{sec:arch-config}

Three IntLLM scales validate the Phase~D scaling-chain
contribution~(C1, Section~\ref{sec:results-scaling}). All three share
the BitLinear forward path of Section~\ref{sec:arch-bitlinear} and the
per-block layout of Section~\ref{sec:arch-block}; they differ only in
hidden dimension and depth.

\begin{table}[t]
\centering
\small
\caption{IntLLM configuration matrix. Parameter counts include the
ternary BitLinear weights, FP32 norms, and the embedding+lm\_head
matrices at full precision in the trained checkpoint. The
\texttt{.fjm}~v9 export format reduces the trained checkpoint to a
ternary-packed kernel binary.}
\label{tab:config}
\begin{tabular}{lrrrr}
\toprule
Config  & $d$  & $L$ & Vocab  & Params \\
\midrule
Mini    & 256  & 6   & 32{,}768 & 21.9~M \\
Base    & 384  & 12  & 32{,}768 & 46.4~M \\
Medium  & 512  & 12  & 32{,}768 & 74.5~M \\
\bottomrule
\end{tabular}
\end{table}

The hidden dimensions $\{256, 384, 512\}$ are deliberately chosen as
powers-of-two-friendly values: $d = 2^k$ supports the Walsh-Hadamard
rotation construction of Section~\ref{sec:negatives} without padding,
$d \cdot 4 / 3$ MLP intermediate sizes pack cleanly into ternary kernel
buffers, and the embedding tables fit in $\leq 2^{15}$ vocabulary
without dynamic sizing.

\subsection{Integer-native activation pipeline}
\label{sec:arch-pipeline}

Beyond BitLinear, three additional surface-area changes make the
forward path integer-native end-to-end:

\paragraph{$\gamma$-less RMSNorm.} The upstream MatMul-Free RMSNorm
multiplies its normalized output by a learnable $\gamma$ tensor; we
absorb this scale into the downstream BitLinear's effective $\beta$
calibration constant during checkpoint export, eliminating a per-block
FP multiplication from the kernel-resident forward path. No
training-time loss-of-expressivity penalty was observed in the Phase~D
scaling chain (Section~\ref{sec:results-scaling}).

\paragraph{Sigmoid and SiLU lookup tables.} The HGRN gating activation
$\sigma$ and the MLP gate activation SiLU are replaced at inference
time by 257-entry linear-interpolation lookup tables defined over
$x \in [-8, +8]$, with maximum absolute error $\leq 2 \cdot 10^{-3}$ and
$\leq 1.6 \cdot 10^{-2}$ respectively. The training-time forward path
uses standard PyTorch \texttt{torch.sigmoid} and
\texttt{torch.nn.functional.silu} for gradient stability; the
training-versus-inference equivalence is verified by an
\texttt{fp16-vs-ternary parity gate} (deferred to Path~A Week~3
calendar; Section~\ref{sec:method-trackb}).

\paragraph{Activation quantization at module boundary.} The activation
quantizer $q_a$ of Equation~\ref{eq:actquant} is applied at every
BitLinear input, producing 8-bit integer activations that are
multiplied by ternary weights and accumulated into 16-bit (training) or
32-bit (inference) integers before dequantization. At no point in the
forward path does the kernel-resident inference code path require a
floating-point matmul, satisfying the kernel-context invariant of
zero-FP-hot-path execution.

\subsection{Comparison with prior architectures}
\label{sec:arch-comparison}

Table~\ref{tab:arch-comparison} positions IntLLM against the closest
prior architectures along the axes that drive the kernel-deployment
constraint (full citations in
\texttt{docs/INTLLM\_RELATED\_WORK.md}~\textsection7).

\begin{table}[t]
\centering
\small
\caption{Architecture comparison on axes relevant to kernel-context
deployment. ``$\checkmark$'' = present, ``$\circ$'' = partial,
``$-$'' = absent. The integer-native column requires
zero floating-point matmuls in the forward path.}
\label{tab:arch-comparison}
\begin{tabular}{l c c c c c}
\toprule
                  & Tern.\ wts & Int.\ acts & No softmax & No FP $\sigma$ & Kernel \\
\midrule
BitNet b1.58 2B4T &$\checkmark$&     $-$    &     $-$    &       $-$      &  $-$  \\
MatMul-Free 2.7B  &$\checkmark$&     $-$    & $\checkmark$&      $-$      &  $-$  \\
TerEffic FPGA     &$\checkmark$&  $\circ$   &     $-$    &       $-$      &  $-$  \\
\textbf{IntLLM}   &$\checkmark$&$\checkmark$&$\checkmark$ &$\checkmark$ &$\checkmark$\\
\bottomrule
\end{tabular}
\end{table}

The composition of the five $\checkmark$'s in IntLLM's row, rather than
any single axis, is the architectural contribution. We do not claim
that any axis is novel in isolation; we claim that the joint
satisfaction of all five is what unlocks the kernel-resident deployment
demonstrated in Section~\ref{sec:results-kernel}.

% ═══════════════════════════════════════════════════════════════════════
\section{Training Methodology and Track~B Interruption Safety}
\label{sec:method}
\label{sec:method-trackb}
\label{sec:method-integrity}

\textit{[STUB — drafted next sub-step (Path A Week 1 sub-step 3).
Covers calibrated gates per scale (Mini~$<$~4.5, Base~$<$~4.2,
Medium~$<$~4.0; calibrated against Phase~D Mini~v2 evidence per
\texttt{FJQ\_PHASE\_D\_GATE\_CALIBRATION.md}); checkpoint atomicity
(\texttt{.tmp}~+~\texttt{os.replace}); step-idle watchdog (1800~s);
HF-streaming retry on transient network errors; bilingual stream
60:40 ID:EN sampling; like-for-like 24K-step ablation budget per
Q5~baseline; research-integrity rules adopted from CLAUDE.md
\textsection6.9~R1-R7 (canonical-protocol benchmarks, literature
review precedes algorithm design, baseline parity, calibrated-once
decompositions, outlier-handling required, results section last,
pre-publication audit gate).]}

% ═══════════════════════════════════════════════════════════════════════
\section{Phase~D Scaling-Chain Results (C1)}
\label{sec:results-scaling}

\textit{[STUB — drafted Week 2. Tables~1+2, Figure~1, narrative.]}

% ═══════════════════════════════════════════════════════════════════════
\section{Bilingual Extension and Q5 Baseline (C2)}
\label{sec:bilingual}

\textit{[STUB — drafted Week 2. Phase~E1 corpus v1.0 description,
Q5 measurement protocol, Table~5 partial, Figure~4.]}

% ═══════════════════════════════════════════════════════════════════════
\section{Honest Negative Results on Outlier-Handling Features (C4)}
\label{sec:negatives}

\textit{[STUB — drafted Week 2. Lifts content from
\texttt{FJQ\_PHASE\_E\_E2\_BILINGUAL\_CALIB\_DECISION.md} v1.0 and
\texttt{FJQ\_PHASE\_E\_E2\_HADAMARD\_DECISION.md} v1.0, including the
3-cause attributions and Phase F.5/F.6 future-work pointers. Tables~5
full and~6, Figure~3.]}

% ═══════════════════════════════════════════════════════════════════════
\section{In-Kernel Deployment via FajarOS Nova (C3)}
\label{sec:results-kernel}

\textit{[STUB — drafted Week 3. FajarOS Nova v3.9.0 IntLLM kernel-path;
Table~4 binary-size + tok/s + watts; Figure~2 binary budget breakdown;
4-invariant regression gate \texttt{make test-intllm-kernel-path}.]}

% ═══════════════════════════════════════════════════════════════════════
\section{Compiler-Verified Safety (C5)}
\label{sec:safety}

\textit{[STUB — drafted Week 3. \texttt{@kernel} / \texttt{@device} /
\texttt{@noinline} annotations; type-system enforcement of the
no-heap, no-tensor, no-raw-pointer invariants at the kernel-LLM
deployment boundary; Table~7 compile-time error examples (KE001
HeapAllocInKernel, KE002 TensorInKernel, etc.); Figure~5 compilation
pipeline diagram with the @kernel analyzer pass.]}

% ═══════════════════════════════════════════════════════════════════════
\section{Discussion and Future Work}
\label{sec:discussion}

\textit{[STUB — drafted Week 3. Phase F.5 PTQ-style post-training
calibration; Phase F.6 post-hoc QuaRot on pre-trained checkpoint;
Phase F.7-F.9 deferred E2.2-E2.5 sub-features; Phase~E3 bilingual
scale-up to Base/Medium; Phase~F~Tier-3 tax-vertical; Phase~D Stretch
training. Each Phase~F entry maps to a sub-task table in
\texttt{FJQ\_PHASE\_F\_TAX\_VERTICAL\_ROADMAP.md} v1.3.]}

% ═══════════════════════════════════════════════════════════════════════
% Bibliography (BibTeX entries to be generated from
% docs/INTLLM_RELATED_WORK.md §10 References + verified against arXiv
% per CLAUDE §6.9 R7 audit gate at first-draft time)
% ═══════════════════════════════════════════════════════════════════════

\bibliographystyle{plainnat}
\bibliography{refs}

\end{document}
